\section{Patient Clustering and Phenotyping Methodologies}\label{sec:rw_clustering}

The unsupervised stratification of patient populations into homogeneous phenotypes is a long-standing problem in healthcare analytics, and clustering is the standard tool used to address it. The two works reviewed below are representative of the dominant line of research: a systematic comparison of clustering algorithms for EHR data, and an applied clustering-based decision-support system. A third work, which couples mixed-type clustering with pathway construction, is the closest pipeline to the one proposed in this thesis and is therefore used to bridge towards the next section.

Aljohani \cite{aljohani2024} addresses the problem of selecting an appropriate clustering algorithm for patient stratification over EHR data. The author frames algorithm selection itself as a multi-criteria decision-making problem and conducts a comparative assessment of eight widely used clustering algorithms, namely K-Means, DBSCAN, hierarchical clustering, Mean Shift, Affinity Propagation, Spectral Clustering, Gaussian Mixture Models, and Self-Organizing Maps.

Methodologically, the comparison is driven by a Fuzzy Technique for Order of Preference by Similarity to Ideal Solution (Fuzzy TOPSIS), aggregating expert judgments over a set of evaluation factors that include cluster quality metrics, scalability, robustness to noise, cluster shape and density, interpretability, and stability. The study synthesizes the relative strengths and weaknesses of each algorithm across these dimensions and provides practitioners with guidance on which algorithm best fits a given patient-stratification setting.

The principal limitation of this work, with respect to the present thesis, is that every one of the eight evaluated algorithms operates natively only on numerical feature spaces. Clinical records, however, are intrinsically mixed: they combine continuous physiological measurements with categorical attributes such as diagnostic codes, drug classes, and gender. Applying these algorithms to clinical data therefore requires encoding categorical variables numerically, which inflates dimensionality and introduces the geometric distortions discussed in Chapter~\ref{ch:analysis}. Moreover, the study is confined to the algorithm-selection phase and does not address any downstream analytical task, such as temporal sequence extraction or the automated interpretation of the resulting clusters.

Yang et al. \cite{yang2024} take an applied perspective, proposing a smart healthcare system that uses clustering on medical data to analyze potential risks and trends within a patient population over a given time frame. The system is explicitly oriented towards hospital management, with the discovered clusters used to compare health profiles across patient subgroups, including a comparison of indigenous and non-indigenous patients in terms of diseases, time, and numbers.

The approach integrates clustering output with evidence-based medicine, examining the similarities and differences between the obtained clusters in order to derive personalized medical services and to inform managerial decisions. In this respect the goal is close in spirit to HospitAI, in that the clustering result is intended to be consumed by non-technical staff rather than left as a purely statistical artifact.

Nevertheless, the bridge from numbers to interpretation remains manual. The clustering operates on homogeneous numerical features, so categorical clinical context must again be discarded or encoded during preprocessing, and the discovered clusters are delivered as static reports whose interpretation depends entirely on expert review. There is no mechanism that automatically translates each cluster into a clinical narrative, nor any facility for an end user to interrogate the results conversationally.

Closest to the architecture proposed in this thesis is the two-step pipeline of Najjar et al. \cite{najjar2018}, which mines patient treatment pathways from administrative healthcare databases. Crucially, this work already employs the K-Prototypes algorithm to cluster heterogeneous medical visits, explicitly acknowledging the mixed numerical-categorical nature of clinical data that motivates the same choice in HospitAI. The clustered visits are then assembled into higher-level treatment pathways, which are themselves clustered to distill typical trajectories, and the approach is validated on a real-world cohort of elderly heart-failure patients.

Even this advanced pipeline, however, leaves two gaps that the present work targets. First, the temporal structure of care is modeled through Hidden Markov Models rather than through frequent sequential pattern mining, so the output is a probabilistic transition model rather than an explicit, support-ranked set of clinical sequences that can be traced back to individual patients. Second, the interpretation of the resulting clusters and pathways is left to domain experts; there is no automated, constraint-based verbalizer that converts the statistical output into clinician-readable summaries. These two gaps motivate, respectively, the sequential pattern mining and the LLM interpretation pillars reviewed in the following sections.
